\scanpage{33}
\mspage{16}
Les applications sont \uline{injectives}.
\deleted{16}

\uline{Démonstration}. Il suffit de le prouver pour la seconde application, et on est ramené au cas où $X$ est \inserted{lui-$\unsure{\overline m}$} un groupe alg. \inserted{$G$}, ext. d'une VA par un tore. Alors l'application est un hom de groupes
\[
(16)\qquad \begin{gathered}G(K)\\\Vert\\H^0(K,G)\end{gathered}\longrightarrow H^1(K,H_1(\overline G))
\]
obtenu ainsi. On considère pour tt entier $n$ la suite exacte Kummérienne
\[
0\longrightarrow {}_nG\longrightarrow G\xrightarrow{n\operatorname{id}_G}G\longrightarrow1
\]
d'où suite exacte de cohomologie, qui donne
\[
0\longrightarrow H^0(K,G)_n\longrightarrow H^1(K,{}_nG)\longrightarrow {}_nH^1(K,G)\longrightarrow0
\]
et passant à la limite, on trouve
\[
0\longrightarrow\varprojlim_n H^0(K,G)_n\longrightarrow\varprojlim_n H^1(K,{}_nG)\longrightarrow\varprojlim_n {}_nH^1(K,G)
\]
Le composé de (16) avec l'hom \deleted{\unsure{composé}} \inserted{canonique}
\[
H^1(K,H_1(\overline G))-\varprojlim_n H^1(K,{}_nG)
\]
compte tenu de
\[
H_1(\overline G)=\varprojlim_n {}_nG
\]
n'est autre que l'hom induit par
\[
H^0(K)\longrightarrow\varprojlim H^0(K,G)_n\hookrightarrow\varprojlim H^1(K,{}_nG)
\]
\margin{32}
